\documentclass[11pt,a4paper]{article}
\usepackage[utf8]{inputenc}
\usepackage[T1]{fontenc}
\usepackage[spanish,es-nodecimaldot]{babel}
\usepackage{amsmath,amssymb}
\usepackage[round]{natbib}
\usepackage{booktabs}
\usepackage{graphicx}
\usepackage[margin=2.5cm]{geometry}
\usepackage{hyperref}
\usepackage{xcolor}
\hypersetup{colorlinks=true,linkcolor=blue!50!black,citecolor=blue!50!black,urlcolor=blue!50!black}
\usepackage{caption}
\captionsetup{font=small,labelfont=bf}

\title{\bfseries La similitud encuentra el hecho, no la versión\\
\large Un sello de orden co-entrenado en la clave resuelve el conflicto de versiones\\
en una memoria persistente dentro del modelo}

\author{Maximiliano Speranza\\
\small Investigador independiente, Buenos Aires, Argentina\\
\small \href{https://orcid.org/0009-0005-0413-8554}{ORCID 0009-0005-0413-8554}\\
\small \texttt{maximiliano.speranza@gmail.com}}

\date{Septiembre de 2026}

\begin{document}
\maketitle

\begin{abstract}
\noindent
Un archivo persistente co-entrenado adentro de un modelo de lenguaje aprende a recuperar el hecho
correcto y después no sabe cuál versión de ese hecho rige. Esa falla está medida con precisión. Con
una sola versión archivada el modelo responde $0{,}9974$, y con dos versiones de la misma clave
compitiendo cae a $0{,}4576$, que es el azar entre el valor viejo y el nuevo. La recuperación queda
intacta y lo que falta es el orden.

Mostramos que la información faltante entra en un \textbf{sello de orden} chico y co-entrenado,
agregado a la clave del archivo, y que lo que hace el trabajo es el ordenamiento y no la capacidad
extra. Sobre cinco semillas, las claves revisadas pasan de $0{,}4570 \pm 0{,}0289$ sin sello a
$\mathbf{0{,}9956 \pm 0{,}0029}$ con sello. El control de falsación, un sello con los mismos
parámetros cuyo valor no guarda relación con el turno real de escritura, da $0{,}4768 \pm 0{,}0389$,
adentro de una desviación estándar del piso. Las quince corridas se ordenan perfecto por condición,
sin un solo solape entre celdas.

Otros tres resultados precisan la afirmación. Primero, con tres versiones por clave, cada una escrita
en su propia secuencia para que no sobreviva ningún rastro de recencia en el contenido, preguntar por
el valor superado pasa de $0{,}2917$, que es el azar entre tres, a $0{,}8316$. Segundo, cuando los
turnos de escritura se sortean al azar y ninguna tabla fija de posición puede resolver la tarea, el
sello igual da $0{,}9644$ en la versión vigente, lo que descarta ese atajo. Tercero, y es lo más útil
en la práctica, un sello \emph{corrupto} es peor que no tener sello, y degrada incluso la
identificación del ítem, que baja de $0{,}9974$ a $0{,}8802$. Una marca de tiempo equivocada no es
información faltante sino información dañina.

También separamos dos capacidades que suelen confundirse. Preferir el valor más reciente y usar el
orden se aprenden en momentos distintos, y la vigente satura a los $3000$ pasos mientras el valor
anterior está todavía en cero exacto. Un banco de pruebas que sólo pregunta por el valor vigente no
puede ver esa distinción, porque el atajo de la recencia ya le da la respuesta correcta.
\end{abstract}

\section{El problema, medido}

Las memorias de atención lineal y de regla delta tienen un estado recurrente de tamaño fijo y
sobrescriben al escribir, así que lo viejo se pierde por construcción \citep{deltanet,gateddeltanet}.
La reparación natural es un archivo externo, y una línea de trabajo creciente co-entrena ese archivo
adentro del modelo en vez de enchufarle un recuperador congelado \citep{colmlm,persistmem,lmlm}.

Esa reparación funciona para recuperar y deja un agujero específico. En nuestra etapa previa, un
archivo co-entrenado que guarda vectores producidos por el propio modelo llega a $0{,}9974$ cuando una
clave tiene una sola versión archivada, y se derrumba a $0{,}4576$ cuando esa misma clave tiene dos
versiones compitiendo. El valor $0{,}4576$ no es ruido, es el azar entre la vieja y la nueva. El
modelo encuentra el hecho correcto y no sabe cuál versión rige.

Es el mismo modo de falla que ya habíamos medido en un índice no paramétrico, donde la geometría
agrupaba las versiones de un recuerdo perfecto y recuperaba la \emph{más vieja}. Dos mecanismos sin
nada en común fallan igual, lo que sugiere que el diagnóstico es una propiedad de la recuperación por
similitud y no de una implementación particular. La similitud identifica el ítem. El orden es
información que la similitud no contiene.

\section{Trabajo relacionado}

\textbf{Híbridos intra-secuencia.} \citet{hola} combina un estado de atención lineal con una memoria
exacta para lo que ese estado olvida, y \citet{ham} junta un estado recurrente con atención que lo
complementa sólo donde el estado predice mal. Los dos son híbridos de un estado compresivo con un
almacén exacto, y en los dos ese almacén vive adentro de un solo paso hacia adelante, así que la
pregunta de cuál de dos versiones escritas en sesiones \emph{distintas} rige no puede aparecer.

\textbf{Memoria persistente co-entrenada.} \citet{colmlm} es lo más cercano en mecanismo, porque emite
consultas continuas desde los estados ocultos del modelo para recuperar de una base de conocimiento,
con $135$M a $360$M de parámetros. \citet{persistmem} acumula un banco de memoria entre turnos sobre
un GPT-2 congelado, de modo que hechos dichos en una sesión se recuerdan después. Ninguno versiona un
hecho, ninguno pone orden en la clave, y ninguno pregunta jamás por un valor superado. Co-LMLM trata
la base como una foto estática, y \citet{persistmem} aplica un decaimiento temporal uniforme, que
atenúa lo viejo en vez de ordenarlo.

\textbf{Interferencia proactiva.} \citet{unabletoforget} es lo más cercano en la \emph{tarea}, porque
actualiza claves varias veces adentro del contexto y pide el valor vigente con las versiones viejas
presentes como interferencia. Documenta el mismo atajo de recencia que encontramos nosotros. Es un
banco de pruebas y un diagnóstico, sus autores declaran que no ofrecen mitigación arquitectónica,
evalúa sólo modelos grandes preentrenados, y nunca pide el valor anterior. Ese último punto nos
importa, porque el atajo de la recencia responde solo la pregunta por el valor vigente.

\textbf{Versionado en capa de aplicación.} Existen sistemas que versionan documentos con metadatos y
punteros explícitos, en cañerías de recuperación montadas sobre un modelo congelado. El sello que
estudiamos no es eso. Vive adentro de la clave de recuperación del propio modelo y se entrena junto al
lector, en vez de estar en un campo de metadatos que consume un recuperador aparte.

\textbf{Marcas de tiempo en la entrada.} \citet{tslm} inyecta marcas temporales en la entrada para
flujos de eventos, pasado, presente y futuro. Eso no es dos versiones del mismo hecho adentro de una
memoria, y la marca no entra en ninguna clave de recuperación.

\section{Montaje experimental}

El sustrato es un modelo de lenguaje chico con un archivo co-entrenado, $64$ dimensiones ocultas, un
archivo de $9$ entradas y una tarea sintética entre secuencias. Los hechos se escriben en una
secuencia, se revisan en otra y se consultan en una tercera, \textbf{con el estado recurrente reseteado
entre secuencias}. La recuperación entonces no se puede resolver arrastrando activaciones, y el piso
es el azar por construcción.

El archivo guarda vectores producidos por el modelo, no tokens. Escribir la entrada $i$ desde el
estado $h_i$ da la clave $k_i = h_i W_k$ y el valor $v_i = h_i W_v$, y leer forma una consulta desde el
estado que consume y atiende sobre el archivo. El \textbf{sello} agrega a la clave un embedding
aprendido del turno de escritura,
\[
k_i \;=\; h_i W_k \;+\; \mathrm{ord}[\,t_i\,],
\]
donde $t_i$ es el turno en que se escribió la entrada $i$ y $\mathrm{ord}$ es una tabla entrenada. Son
$384$ parámetros efectivos. Y algo decisivo, \textbf{el archivo se baraja por muestra} y cada sello
viaja pegado a su propia entrada. Sin barajar, el turno de escritura coincidiría siempre con la
posición en el tensor y el experimento mediría mirar el final del archivo, que es posición y no orden.

Se comparan tres condiciones en todo el trabajo. \texttt{ninguno} no tiene sello. \texttt{sello} es lo
de arriba. \texttt{barajado} es el control de falsación, con una tabla del mismo tamaño y entrenada
igual, cuyo valor de sello \emph{no guarda relación} con el turno real de escritura.

\section{Resultado 1. El sello levanta el conflicto de versiones, y es el orden lo que lo levanta}

Cinco semillas, dos versiones por clave revisada, $4000$ pasos.

\begin{center}
\begin{tabular}{lccc}
\toprule
condición & global & \textbf{claves revisadas} & claves no revisadas \\
\midrule
\texttt{ninguno} & $0{,}7273 \pm 0{,}0146$ & $0{,}4570 \pm 0{,}0289$ & $0{,}9977$ \\
\textbf{\texttt{sello}} & $\mathbf{0{,}9970 \pm 0{,}0020}$ & $\mathbf{0{,}9956 \pm 0{,}0029}$ & $0{,}9984$ \\
\texttt{barajado} & $0{,}7374 \pm 0{,}0194$ & $0{,}4768 \pm 0{,}0389$ & $0{,}9979$ \\
\bottomrule
\end{tabular}
\end{center}

La ganancia sobre no tener sello es de $+0{,}5385$ y sobre el control barajado de $+0{,}5187$.
\textbf{Las cinco semillas con sello caen entre $0{,}9909$ y $0{,}9987$, y la mejor semilla de las
otras dos condiciones llega a $0{,}5156$. Las quince corridas se ordenan perfecto por condición, sin
un solo solape.} Con esa distancia, el intervalo de confianza es un trámite.

\textbf{Lo que hace válido el resultado es la celda que no ganó.} La condición \texttt{barajado} tiene
exactamente los mismos parámetros que \texttt{sello}, entrenados igual, y lo único que cambia es que el
sello no guarda relación con el turno real. Sin esa celda, el $+0{,}5385$ se explicaría igual de bien
diciendo que se le agregaron parámetros al lector. Con ella, la explicación por capacidad queda
descartada de forma directa.

\textbf{La aritmética es simétrica.} El conflicto de versiones costaba $-0{,}5398$ y el sello devuelve
$+0{,}5385$. Repone lo que la similitud había perdido y nada más, porque las claves no revisadas se
quedan en $0{,}998$, donde ya estaban.

\section{Resultado 2. Tres versiones, y preguntar por el valor superado}

Un primer intento con tres versiones se rompió solo, y cómo se rompió vale registrarlo. Al escribir las
dos revisiones adentro de la misma secuencia, el estado en $v_3$ arrastraba un rastro de $v_2$, así que
el orden quedaba escrito en el \emph{contenido} sin que ningún metadato lo pusiera. A $4000$ pasos era
invisible porque ninguna condición resolvía la pregunta difícil, y a $12000$ pasos la condición sin
sello llegó a $0{,}97$ y delató la fuga.

Con cada versión escrita en su propia secuencia y el estado reseteado entre ellas, $v_2$ y $v_3$ son
cada una primera y única mención de su secuencia, indistinguibles en recencia. Lo único en el sistema
que dice cuál vino antes es el sello. Tres semillas, $12000$ pasos.

\begin{center}
\begin{tabular}{lccc}
\toprule
condición & vigente & \textbf{anterior} & claves de una sola versión \\
\midrule
\texttt{ninguno} & $0{,}3090$ & $0{,}2917 \pm 0{,}0207$ & $0{,}9436$ \\
\textbf{\texttt{sello}} & $0{,}9774$ & $\mathbf{0{,}8316 \pm 0{,}2399}$ & $0{,}9974$ \\
\texttt{barajado} & $0{,}2682$ & $0{,}2804 \pm 0{,}0159$ & $0{,}8802$ \\
\bottomrule
\end{tabular}
\end{center}

Sin sello el modelo se queda en $1/3$, el azar entre tres versiones, y al mismo tiempo identifica bien
el ítem cuando no hay conflicto, con $0{,}9436$. La fuga quedó tapada, y la distancia respecto del
$0{,}97$ del diseño roto mide exactamente cuánta información de orden se estaba colando por el
contenido.

\textbf{Este número depende del presupuesto y lo informamos así.} Con cinco semillas el valor anterior
baja a $0{,}7667$ con desviación $0{,}2340$, y sus valores son $0{,}9766$, $0{,}9635$, $0{,}5547$,
$0{,}8594$ y $0{,}4792$. El mecanismo no falla, sino que dos de cinco semillas no convergieron a
$12000$ pasos. Se probó de forma directa, y la semilla 2 extendida a $24000$ pasos pasó de $0{,}5547$
a $0{,}9531$. El veredicto honesto es condicional, y dice que el sello permite responder por la versión
anterior con un presupuesto que $12000$ pasos no siempre alcanza. \textbf{En este régimen el tiempo de
convergencia es parte del fenómeno y no un detalle de implementación.}

\section{Resultado 3. No es una tabla de posiciones}

En el diseño anterior los turnos de escritura eran fijos por rol, así que el modelo podía aprender que
los embeddings $9$, $10$ y $11$ significan versión vigente. Eso es una correspondencia entre posición y
rol, no un orden, y el control barajado no las separa, porque aleatorizar el sello rompe la tabla igual
que rompe el orden.

Entonces sorteamos doce turnos al azar de un rango de treinta y dos por muestra, los ordenamos y los
asignamos respetando el orden real de escritura. El turno $9$ puede ser una primera versión en una
muestra y una tercera en la siguiente, de modo que ninguna tabla fija resuelve la tarea. Tres semillas.

\begin{center}
\begin{tabular}{lcc}
\toprule
condición & \textbf{vigente} & anterior \\
\midrule
\texttt{ninguno} & $0{,}3247$ & $0{,}2977 \pm 0{,}0326$ \\
\textbf{\texttt{sello}} & $\mathbf{0{,}9644}$ & $0{,}3142 \pm 0{,}5330$ \\
\texttt{barajado} & $0{,}2804$ & $0{,}2925 \pm 0{,}0266$ \\
\bottomrule
\end{tabular}
\end{center}

El resultado en la versión vigente sobrevive a los turnos móviles, lo que \textbf{descarta el atajo de
la tabla y valida hacia atrás el Resultado 1}, porque si el mecanismo hubiera sido una tabla, mover los
turnos habría roto también la vigente.

La columna del valor anterior es \textbf{bimodal, no ruidosa}, con valores $0{,}0052$, $0{,}9297$ y
$0{,}0078$. Dos semillas caen en un atajo de recencia y responden \emph{por debajo} del azar, porque
devuelven la versión vigente cuando se pide la anterior, y una semilla aprende la comparación completa.
Informamos la distribución y no la media, porque el promedio de $0{,}31$ no describe a ninguna de las
tres corridas. El resumen honesto es que el lector aprende a leer el reloj para saber qué hora es
ahora, no para reconstruir la secuencia de lo que pasó.

\section{Resultado 4. Un sello mentiroso es peor que ninguno}

Leer la fila \texttt{barajado} del Resultado 2 como un resultado en sí mismo, y no sólo como control,
da el hallazgo con más filo práctico. Un sello corrupto deja la versión vigente en $0{,}2682$, por
\emph{debajo} del $0{,}3090$ de no tener sello, y degrada la identificación del ítem de $0{,}9974$ a
$0{,}8802$, que es una capacidad que no tiene nada que ver con ordenar.

\textbf{El lector se apoya en el sello.} Una vez que lo hace, una marca de tiempo corrupta no es
información faltante sino información dañina, y el daño se contagia a la calidad de la recuperación.
Para cualquier sistema que le ponga marca de tiempo a sus entradas de memoria, esto dice que un sello
viejo o equivocado es peor que un campo ausente.

\section{Dos capacidades que se aprenden en momentos distintos}

Preferir el valor más reciente y usar el orden son cosas distintas y separables, y el orden entre las
dos es estable entre réplicas aunque el momento no lo sea. En las dos réplicas con tasa de aprendizaje
$10^{-3}$ la versión vigente satura a los $3000$ pasos, con $0{,}9583$ en ambas, mientras el valor
anterior está todavía en $0{,}0000$. Después el valor anterior despega, y cuánto más tarde varía. Una
réplica llega a $0{,}2500$ a los $4000$ pasos, $0{,}6250$ a los $5000$ y $0{,}9583$ a los $6000$,
mientras que la otra a los $4000$ está en $0{,}0208$ y no pasa $0{,}9$ hasta los $8000$. Informamos las
dos porque el promedio no describiría a ninguna. Lo que replica es el \textbf{orden}, ya que en toda
corrida la versión vigente se resuelve miles de pasos antes que la anterior, y en el paso donde la
vigente ya está en $0{,}9583$ la anterior está en cero exacto.

Esto tiene una consecuencia directa para la evaluación. \textbf{Un banco de pruebas que sólo pregunta
por el valor vigente no puede ver la distinción}, porque el atajo de la recencia da la respuesta
correcta. Es precisamente lo que \citet{unabletoforget} no hace, y por eso consideramos que la pregunta
por el valor anterior es la sonda más informativa.

\section{Límites}

\textbf{Es suficiencia, no inferencia.} En esta tarea la versión vigente es siempre la del turno más
alto, así que la política aprendible es quedarse con el sello más grande. El resultado dice que la
información que a la similitud le falta \emph{entra en un sello} y que el lector la sabe usar. No dice
que el modelo razone sobre el tiempo.

\textbf{El sello se le da, no se infiere.} Mostramos que el lector puede usar el metadato, no que pueda
fabricarlo. No es un oráculo irreal, porque el turno de escritura es gratis para cualquier índice, a
diferencia de saber cuál versión es la correcta, pero acota la afirmación.

\textbf{Escala.} Un modelo de $64$ dimensiones, un archivo de $9$ entradas y una tarea sintética. La
dificultad medida es de alineación y de conflicto de versiones, no de discriminación entre muchos
candidatos.

\textbf{Un archivo fresco.} En estos experimentos el archivo se recomputa en cada paso. Un trabajo
posterior donde las entradas envejecen muestra que el daño es gradual y no un acantilado, y las claves
revisadas bajan de $0{,}9987$ a $0{,}9115$ a medida que el coseno entre el marco de escritura y el de
lectura cae a $0{,}78$.

\section{Conclusión}

La similitud identifica el ítem, y el orden es información que la similitud no contiene y que hay que
llevar aparte. En una memoria persistente co-entrenada adentro de un modelo de lenguaje chico,
llevarla como un sello en la clave lleva el conflicto de versiones del azar a casi el techo, y un
control con los mismos parámetros y un sello sin significado muestra que lo que trabaja es el
ordenamiento y no la capacidad. El mismo mecanismo vuelve respondible la pregunta por el valor
superado, que ningún banco de pruebas del área hace hoy, y viene con una advertencia que se generaliza
más allá de nuestro montaje. Una vez que un lector aprende a apoyarse en un sello, un sello equivocado
cuesta más que no tener ninguno.

\section*{Reproducibilidad y disponibilidad de datos}

Cada experimento acá informado fue pre-registrado con sus predicciones y umbrales hasheados antes de la
corrida correspondiente. Las celdas que no cumplieron sus criterios se informan como incumplimientos,
incluido el valor anterior dependiente del presupuesto del Resultado 2 y la columna bimodal del
Resultado 3. Todos los pre-registros con sus hashes SHA, los scripts de análisis, los resultados crudos
por semilla y los informes generados por máquina están archivados en el repositorio del proyecto,
\url{https://github.com/SperanzaMax/delta-archivo}, cuyo historial de commits lleva las marcas de
tiempo que establecen el orden entre pre-registro y resultado.

\section*{Declaraciones}

\textbf{Financiamiento.} Ninguno. \textbf{Intereses en competencia.} Ninguno. \textbf{Aprobación
ética.} No aplica, no hubo sujetos humanos ni animales. \textbf{Disponibilidad de datos.} Ver arriba.
\textbf{Contribución de autoría.} Autor único.

\begin{thebibliography}{99}
\bibitem[Cui(2026)]{hola} Cui, W. (2026). A Hippocampus for Linear Attention. An Exact Memory for What the Recurrent State Forgets. arXiv:2607.02303.
\bibitem[Lufkin et al.(2026)]{ham} Lufkin, L., Figliolia, T., Millidge, B., y Krishnamurthy, K. (2026). Hybrid Associative Memories. arXiv:2603.22325.
\bibitem[Feldman et al.(2026)]{colmlm} Feldman, Y., Zhao, L., Godey, N., Go, D., Hua, Y., Weinberger, K. Q., Sun, J. J., y Artzi, Y. (2026). Co-LMLM. Continuous-Query Limited Memory Language Models. arXiv:2607.07707.
\bibitem[Jeong(2026)]{persistmem} Jeong, H. (2026). Trained Persistent Memory for Frozen Decoder-Only LLMs. arXiv:2603.22329.
\bibitem[Zhao et al.(2025)]{lmlm} Zhao, L., Zalouk, S., Belardi, C. K., Lovelace, J., Zhou, J. P., Noonan, R. T., Go, D., Weinberger, K. Q., Artzi, Y., y Sun, J. J. (2025). Pre-training Limited Memory Language Models with Internal and External Knowledge. arXiv:2505.15962.
\bibitem[Wang y Sun(2025)]{unabletoforget} Wang, C., y Sun, J. V. (2025). Unable to Forget. Proactive Interference Reveals Working Memory Limits in LLMs Beyond Context Length. arXiv:2506.08184.
\bibitem[Schlag et al.(2021)]{deltanet} Schlag, I., Irie, K., y Schmidhuber, J. (2021). Linear Transformers Are Secretly Fast Weight Programmers. ICML 2021.
\bibitem[Yang et al.(2025)]{gateddeltanet} Yang, S., Kautz, J., y Hatamizadeh, A. (2025). Gated Delta Networks. Improving Mamba2 with Delta Rule. ICLR 2025. arXiv:2412.06464.
\bibitem[Rajaby Faghihi y Kordjamshidi(2021)]{tslm} Rajaby Faghihi, H., y Kordjamshidi, P. (2021). Time-Stamped Language Model. Teaching Language Models to Understand the Flow of Events. NAACL 2021. arXiv:2104.07635.
\end{thebibliography}

\end{document}
